\section{Congruent coefficient algebras, after Maranda}
\label{sec:congruent-coefficients}

Sufficiently close coefficient algebras have the same lattice
Grothendieck group for a fixed finite group.  This transfers the
truncated-polynomial obstruction to equations with small nonzero
coefficients, including totally ramified field orders.
The mechanism is Maranda's averaging and successive approximation:
\citet[Theorem~2, pp.~347--348]{Maranda1953} proves rigidity of
integral representations from sufficiently close reductions, and
his Theorem~3, pp.~348--351, lifts modular subconstituents with the
precision improvement used below.  We give the coefficient-algebra
comparison and its exact-sequence compatibility explicitly, since
these are the assertions needed for the present Grothendieck groups.

\begin{theorem}\label{thm:congruent-coefficients}
Let $A$ and $B$ be nonzero commutative algebras finite free over
$\Z_p$, let $G$ be a finite group, and put $v=v_p(|G|)$.
For an integer $r>2v$, fix a unital algebra isomorphism
\[
 \alpha:A/p^rA\xrightarrow{\sim}B/p^rB.
\]
It determines a canonical rank-preserving bijection between the
isomorphism classes in $\Lat_A(G)$ and $\Lat_B(G)$.  This bijection
preserves direct sums and the triples of classes represented by
coefficient-split short exact sequences.  Hence it induces an
isomorphism
\[
 K_0(\Lat_A(G))\xrightarrow{\sim}K_0(\Lat_B(G)).
\]
The isomorphism respects tensor and exterior powers and representations
defined over $\Z_p$.  It commutes with restriction and induction for
subgroups of $G$.  The assertion concerns isomorphism classes and
exact-sequence relations; no equivalence of the categories of all
morphisms is asserted.
\end{theorem}

\subsection{Correction and rigidity}

We first establish the two estimates that make the comparison
independent of choices.  Let $R$ be a finite free $\Z_p$-algebra,
write $m=|G|$, and suppose that matrices $\rho_g\in\mathrm{GL}_d(R)$
satisfy $\rho_1=I$ and
\[
 D(g,h)=\rho_g\rho_h-\rho_{gh}\in p^s\mathrm{M}_d(R),
 \qquad s>2v.
\]
The matrices need not yet define an action.  Set
\[
 a_g=\frac1m\sum_{k\in G}D(g,k)\rho_k^{-1},\qquad
 \rho'_g=\rho_g-a_g.
\]
Division by $m$ is legitimate here: every numerator entry is
divisible by $p^s$, and $m/p^v$ is a unit.  Thus
$a_g\in p^{s-v}\mathrm{M}_d(R)$.  The new matrices are invertible
and $\rho'_1=I$.

Associativity gives the exact identity
\[
 \rho_gD(h,k)-D(gh,k)+D(g,hk)-D(g,h)\rho_k=0.
\]
Multiply on the right by $\rho_k^{-1}$ and average.  If
\[
 E(g,h)=a_g\rho_h-\frac1m\sum_kD(g,hk)\rho_k^{-1},
\]
reindexing the first sum shows that
\begin{align*}
 E(g,h)
 &=\frac1m\sum_kD(g,hk)
       (\rho_{hk}^{-1}\rho_h-\rho_k^{-1})\\
 &=\frac1m\sum_kD(g,hk)\rho_{hk}^{-1}D(h,k)\rho_k^{-1}.
\end{align*}
Consequently $E(g,h)\in p^{2s-v}\mathrm{M}_d(R)$, and the new
defect satisfies
\begin{equation}\label{eq:congruence-improvement}
 \rho'_g\rho'_h-\rho'_{gh}
   =-E(g,h)+a_ga_h\in p^{2s-2v}\mathrm{M}_d(R).
\end{equation}
This calculation uses no cohomology action associated with an
approximate representation.

Iteration replaces $s$ by
$s_j=2v+2^j(s-2v)$, which tends to infinity.  Completeness and
separatedness of $R$ give an exact action $\sigma$ satisfying
\begin{equation}\label{eq:congruence-corrected}
 \sigma_g\equiv\rho_g\pmod {p^{s-v}}.
\end{equation}
The construction preserves block diagonal and block upper triangular
matrix algebras.  In either case its diagonal blocks are exactly
the separate constructions on the original diagonal blocks.

For rigidity, let $\sigma,\tau$ be exact actions, and suppose an
invertible matrix $C$ satisfies
\[
 \tau_gC-C\sigma_g\in p^s\mathrm{M}_d(R),\qquad s>v.
\]
The averaged intertwiner
\begin{equation}\label{eq:congruence-intertwiner}
 F=\frac1m\sum_{g\in G}\tau_gC\sigma_g^{-1}
\end{equation}
is integral and congruent to $C$ modulo $p^{s-v}$.  It is therefore
invertible.  Reindexing the sum proves $\tau_hF=F\sigma_h$ for
every $h\in G$.  In particular exact actions congruent modulo
$p^s$, for $s>v$, are isomorphic.

\subsection{Transport of objects and every exact relation}

Lift $\alpha$ to a unital $\Z_p$-linear isomorphism
$\theta:A\to B$.  Such a lift exists: $1$ is primitive in each
finite free $\Z_p$-module, since $1\in pA$ would imply $A=pA$,
contrary to Nakayama's lemma.  Extend $1$ to bases and lift the
matrix of $\alpha$, fixing its first column.  Its determinant is a
unit.  The inverse map lifts $\alpha^{-1}$, and
\begin{equation}\label{eq:congruence-multiplication}
 \theta(ab)-\theta(a)\theta(b)\in p^rB.
\end{equation}

Choose an $A$-basis for a lattice and apply $\theta$ entrywise to
its action matrices $\rho_g$.  The resulting matrices are
invertible, since their products with the entrywise images of
$\rho_g^{-1}$ are congruent to $I$ modulo $p^r$.
Their multiplication defect lies in $p^r$.  The correction above
produces an exact $B$-action with
\[
 \sigma_g\equiv\theta(\rho_g)\pmod {p^{r-v}}.
\]
Since $r-v>v$, rigidity makes its isomorphism class unique among
actions with this congruence.  Two lifts of $\alpha$ differ by
$p^r$, so their corrected actions are isomorphic.  For a change
of basis $C$, the invertible matrix $\theta(C)$ is an intertwiner
modulo $p^{r-v}$ by \eqref{eq:congruence-multiplication};
\eqref{eq:congruence-intertwiner} makes it an exact intertwiner.
Thus the transported class depends only on $\alpha$ and the
original isomorphism class.

Transporting back with $\theta^{-1}$ gives an exact action
congruent to $\rho$ modulo $p^{r-v}$: the first correction and the
second correction both have this divisibility.  Rigidity identifies
the round trip with the original action.  The two transports are
therefore inverse bijections.

Now choose an $A$-splitting of an admissible short exact sequence.
In its adapted basis the action matrices are block upper triangular.
Entrywise transport and every correction preserve that form, and
their diagonal blocks correct the actions on the subobject and
quotient separately.  The corrected matrices thus give an admissible
short exact sequence with the three transported classes.
The inverse transport has the same property.  Block diagonal
matrices prove direct-sum compatibility.  Passing from the free
abelian groups on isomorphism classes to their quotients by all
exact-sequence relations proves the claimed $K_0$ isomorphism.

Tensor and exterior power matrices are integral polynomials in
the entries.  Applying $\theta$ before or after those polynomials
changes the result by $p^r$.  The corrected exact actions are
therefore congruent modulo $p^{r-v}$ and are isomorphic by rigidity.
Matrices over $\Z_p$ are fixed throughout.  For a subgroup
$H\le G$, put $v_H=v_p(|H|)\le v$.  Restriction followed by
correction, or correction followed by restriction, gives exact
$H$-actions congruent modulo $p^{r-v}$, which exceeds $v_H$.
For induction, fixed coset representatives give matrices with blocks
equal to action matrices or zero.  The two induced $G$-actions are
congruent modulo $p^{r-v}$, which exceeds $v$.
Rigidity proves both compatibilities and completes the theorem.

